\section*{Summary in plain words}
\addcontentsline{toc}{section}{Summary in plain words}

\paragraph{The question.}
On a Yukon leaderboard, anyone can submit an improvement to a shared problem, and most people start from someone else's earlier submission.
Each submission says which AI model was used, for example a Claude model from Anthropic or a GPT model from OpenAI.
The question is simple: when you build on someone else's work, does it help if their work was made with a different company's model than yours?

\paragraph{What I did.}
I took a frozen copy of the public record of \StudyK{} leaderboards, worked out for every submission which earlier submission it built on, and recorded the company behind each model.
A submission counts as a success if it was promoted, meaning it beat the best score so far and became the new starting point for everyone.
The rules for the analysis were written down and timestamped before the outcomes were joined to the data, so the method could not be tuned to get a particular answer.
A few rough looks at ECDSA.Fail before that point are disclosed in Section~\ref{sec:peek}.

\paragraph{What I found.}
Building on another company's model was the common case.
Among the \PrimaryRowsTotal{} submissions in the main analysis, all building on someone else's work, \RawCrossRows{} built on another company's model and \RawSameRows{} on the same company's.
Their raw success rates were \RawCrossRate{} and \RawSameRate{}.
But that gap disappears when each participant is compared only with themselves, at the same stage of the leaderboard: the chance of success is then almost exactly the same either way.
The best estimate is about \PlainPointPct{} higher odds for cross-company building, and the data rule out an advantage or disadvantage bigger than roughly \PlainHighPct{} up or \PlainLowPct{} down.
All nine alternative ways of running the analysis gave the same answer.

\paragraph{What this does not say.}
It is an average over ten leaderboards, and a single leaderboard can still differ.
It covers only work that people submitted, not experiments they never posted.
It relies on the model each participant declared, which is usually but not always accurate.
And it is about models, not people: other participants bringing fresh ideas can matter even if the brand of model does not.

\paragraph{How to read the numbers.}
Table~\ref{tab:glossary} explains the terms the rest of the paper uses.
The main result in statistical form is an odds ratio of \PooledOR{}, where 1 means no difference, with a 90\% interval of \PooledORNinetyLow{} to \PooledORNinetyHigh{}.

\begin{table}[h]
\centering
\small
\caption{Terms used in this paper.}
\label{tab:glossary}
\begin{tabularx}{\linewidth}{lX}
\toprule
Term & Meaning \\
\midrule
Submission, row & One entry on a leaderboard: code plus a public note. \\
Promotion & The submission beat the best score so far and became the new default starting point. This is the success measure. \\
Parent & The earlier submission that a submission built on. \\
Family & The company behind the declared model, such as Anthropic, OpenAI, Google or xAI. \\
Cross-family step & A submission whose model comes from a different company than its parent's model. \\
Solver & A participant on the leaderboard. \\
Odds ratio & How the odds of promotion compare between two groups. 1 means equal, 1.5 means 50\% higher odds, 0.67 means a third lower. \\
Interval & The range of values the data are consistent with. A 90\% interval is the range used to decide whether an effect is negligible. \\
Margin & The range of odds ratios, fixed in advance, that counts as ``no meaningful effect'': 0.67 to 1.50 here. \\
Pre-registration & The written analysis plan, locked and timestamped before the outcomes were joined to the data. \\
\bottomrule
\end{tabularx}
\end{table}
